\par\textbf{Procedimiento completo para comparar tu solución.} El caudal capturado es eficiencia de captura por contenido líquido por flujo de volumen interceptado. El balance térmico suma calor sensible del agua y convección; la eficacia de entrega divide la demanda eléctrica.

\begin{enumerate}

\item Convertir contenido líquido de g/m3 a kg/m3.

\[\fitmath{\mathrm{mw}=0.7\cdot 0.0004\cdot 80\cdot 0.12\approx 0.002688\,\mathrm{kg/s}}\]

\item Diferencia térmica igual en kelvin y Celsius.

\[\fitmath{\mathrm{dT}=5--10\approx 15\,\mathrm{K}}\]

\item Calor sensible del agua capturada.

\[\fitmath{\mathrm{Pw}=\mathrm{mw}\cdot 4180\cdot \mathrm{dT}}\]

\[\fitmath{\mathrm{Pw}=(0.002688)\cdot 4180\cdot (15)\approx 168.5376\,\mathrm{W}}\]

\item Pérdida convectiva de la superficie.

\[\fitmath{\mathrm{Pc}=80\cdot 0.5\cdot \mathrm{dT}}\]

\[\fitmath{\mathrm{Pc}=80\cdot 0.5\cdot (15)\approx 600\,\mathrm{W}}\]

\item Entrada eléctrica con eficacia menor que uno.

\[\fitmath{\mathrm{Pe}=\frac{\mathrm{Pw}+\mathrm{Pc}}{0.85}}\]

\[\fitmath{\mathrm{Pe}=\frac{(168.5376)+(600)}{0.85}\approx 904.1619\,\mathrm{W}}\]

\item Comparación aislada con fusión de hielo a cero grados; no es el proceso principal.

\[\fitmath{\mathrm{Platent}=\mathrm{mw}\cdot 334000}\]

\[\fitmath{\mathrm{Platent}=(0.002688)\cdot 334000\approx 897.792\,\mathrm{W}}\]

\end{enumerate}
